\subsection{Erstellen von Dokumenten ohne passende Dokumentenklasse}
\label{latex-ohne-dokumentenklasse}
Für Dokumente, für die es noch keine eigene Dokumentenklasse gibt (z.\,B. Aushänge oder Schreiben), wird das \hyperref[latex-fsrbase]{Basispaket} direkt eingebunden:
\begin{enumerate}
	\item Die passende \KOMAScript{}-Klasse aus \autoref{latex-koma} mit \mintinline{latex}{\documentclass} laden, z.\,B. \texttt{scrartcl} für einen Aushang.
	\item Direkt darunter das Basispaket mit \mintinline{latex}{\usepackage{fsrbase}} einbinden.
	\item Den Inhalt zwischen \mintinline{latex}{\begin{document}} und \mintinline{latex}{\end{document}} schreiben.
	\item Das Dokument \hyperref[latex-kompilieren]{kompilieren}.
\end{enumerate}

\begin{Codeblock}[title={Minimales Dokument mit dem Basispaket}]{latex}
\documentclass{scrartcl}
\usepackage{fsrbase}

\begin{document}
\fsrlogo[4cm]

\section*{Überschrift}
Hier steht der Inhalt des Dokuments.
\end{document}
\end{Codeblock}
